\documentclass{amsart}
% For dealing with references we use the comment environment
\usepackage{verbatim}
\newenvironment{reference}{\comment}{\endcomment}
%\newenvironment{reference}{}{}
\newenvironment{slogan}{\comment}{\endcomment}
\newenvironment{history}{\comment}{\endcomment}

% For commutative diagrams we use Xy-pic
\usepackage[all]{xy}

% We use 2cell for 2-commutative diagrams.
\xyoption{2cell}
\UseAllTwocells

% We use multicol for the list of chapters between chapters
\usepackage{multicol}

% This is generally recommended for better output
\usepackage{lmodern}
\usepackage[T1]{fontenc}

% For cross-file-references

% Package for hypertext links:
\usepackage{hyperref}

% For any local file, say "hello.tex" you want to link to please

% Theorem environments.
%
\theoremstyle{plain}
\newtheorem{theorem}[subsection]{Theorem}
\newtheorem{proposition}[subsection]{Proposition}
\newtheorem{lemma}[subsection]{Lemma}

\theoremstyle{definition}
\newtheorem{definition}[subsection]{Definition}
\newtheorem{example}[subsection]{Example}
\newtheorem{exercise}[subsection]{Exercise}
\newtheorem{situation}[subsection]{Situation}

\theoremstyle{remark}
\newtheorem{remark}[subsection]{Remark}
\newtheorem{remarks}[subsection]{Remarks}

\numberwithin{equation}{subsection}

% Macros
%
\def\lim{\mathop{\mathrm{lim}}\nolimits}
\def\colim{\mathop{\mathrm{colim}}\nolimits}
\def\Spec{\mathop{\mathrm{Spec}}}
\def\Hom{\mathop{\mathrm{Hom}}\nolimits}
\def\Ext{\mathop{\mathrm{Ext}}\nolimits}
\def\SheafHom{\mathop{\mathcal{H}\!\mathit{om}}\nolimits}
\def\SheafExt{\mathop{\mathcal{E}\!\mathit{xt}}\nolimits}
\def\Sch{\mathit{Sch}}
\def\Mor{\mathop{\mathrm{Mor}}\nolimits}
\def\Ob{\mathop{\mathrm{Ob}}\nolimits}
\def\Sh{\mathop{\mathit{Sh}}\nolimits}
\def\NL{\mathop{N\!L}\nolimits}
\def\CH{\mathop{\mathrm{CH}}\nolimits}
\def\proetale{{pro\text{-}\acute{e}tale}}
\def\etale{{\acute{e}tale}}
\def\QCoh{\mathit{QCoh}}
\def\Ker{\mathop{\mathrm{Ker}}}
\def\Im{\mathop{\mathrm{Im}}}
\def\Coker{\mathop{\mathrm{Coker}}}
\def\Coim{\mathop{\mathrm{Coim}}}

% Boxtimes
%
\DeclareMathSymbol{\boxtimes}{\mathbin}{AMSa}{"02}

%
% Macros for moduli stacks/spaces
%
\def\QCohstack{\mathcal{QC}\!\mathit{oh}}
\def\Cohstack{\mathcal{C}\!\mathit{oh}}
\def\Spacesstack{\mathcal{S}\!\mathit{paces}}
\def\Quotfunctor{\mathrm{Quot}}
\def\Hilbfunctor{\mathrm{Hilb}}
\def\Curvesstack{\mathcal{C}\!\mathit{urves}}
\def\Polarizedstack{\mathcal{P}\!\mathit{olarized}}
\def\Complexesstack{\mathcal{C}\!\mathit{omplexes}}
% \Pic is the operator that assigns to X its picard group, usage \Pic(X)
% \Picardstack_{X/B} denotes the Picard stack of X over B
% \Picardfunctor_{X/B} denotes the Picard functor of X over B
\def\Pic{\mathop{\mathrm{Pic}}\nolimits}
\def\Picardstack{\mathcal{P}\!\mathit{ic}}
\def\Picardfunctor{\mathrm{Pic}}
\def\Deformationcategory{\mathcal{D}\!\mathit{ef}}

\setlength{\emergencystretch}{3em}
\begin{document}
\title[Stacks Project, Tag 00R7]{1628 --- Nonzero finitely presented base-flat modules with flat closed fibre force intermediate flatness along local maps of essentially finitely presented R-algebras}
\maketitle
\noindent Copyright (C) 2005--2025 Johan de Jong. Stacks Project authors. Source: \href{https://stacks.math.columbia.edu/tag/00R7}{Tag 00R7}.

\noindent Editorial difficulty note (not author text). Difficulty H2: The proof synchronizes two local essentially finite-presentation systems and a module presentation, forces both base and closed-fibre flatness at one finite Noetherian stage by killing finite Tor kernels, then uses faithful flatness to recover intermediate-ring flatness. The non-Noetherian two-level fibre criterion requires a genuinely substantial approximation/homological argument rather than a direct application of the Noetherian criterion..

\noindent Editorial provenance note (not author text). History: excerpt from revision \texttt{a0444\allowbreak 6e57e\allowbreak c1fbc\allowbreak 252a8\allowbreak 71afc\allowbreak ec775\allowbreak 2fb28\allowbreak 07b14}, algebra.tex, lines 32991--33111. Reference presentation was adapted; original-author context and core support blocks were retained or added. The mathematical wording is unchanged. Licensed under GNU FDL 1.2 or later; no invariant sections or cover texts. See ../licenses/COPYING. Context blocks reproduce author definitions and supporting results; they are not additional counted problems.

\begin{lemma}
\label{lemma-limit-no-condition-local}
Suppose $R \to S$ is a local homomorphism of local rings.
There exists a directed set $(\Lambda, \leq)$, and
a system of local homomorphisms $R_\lambda \to S_\lambda$
of local rings such that
\begin{enumerate}
\item The colimit of the system $R_\lambda \to S_\lambda$
is equal to $R \to S$.
\item Each $R_\lambda$ is essentially of finite type
over $\mathbf{Z}$.
\item Each $S_\lambda$ is essentially of finite type
over $R_\lambda$.
\end{enumerate}
\end{lemma}

\begin{proof}
Denote $\varphi : R  \to S$ the ring map.
Let $\mathfrak m \subset R$ be the maximal ideal
of $R$ and let $\mathfrak n \subset S$ be the maximal
ideal of $S$. Let
$$
\Lambda = \{
(A, B)
\mid
A \subset R, B \subset S, \# A < \infty, \# B < \infty, \varphi(A) \subset B
\}.
$$
As partial ordering we take the inclusion relation. For each
$\lambda = (A, B) \in \Lambda$ we let $R'_\lambda$ be
the sub $\mathbf{Z}$-algebra generated by
$a \in A$, and we let $S'_\lambda$ be the sub
$\mathbf{Z}$-algebra generated by $b$, $b \in B$.
Let $R_\lambda$ be the localization of $R'_\lambda$
at the prime ideal $R'_\lambda \cap \mathfrak m$ and let
$S_\lambda$ be the localization of $S'_\lambda$ at
the prime ideal $S'_\lambda \cap \mathfrak n$.
In a picture
$$
\xymatrix{
B \ar[r] &
S'_\lambda \ar[r] &
S_\lambda \ar[r] &
S \\
A \ar[r] \ar[u] &
R'_\lambda \ar[r] \ar[u] &
R_\lambda \ar[r] \ar[u] &
R \ar[u]
}.
$$
The transition maps are clear. We leave the proofs of the other
assertions to the reader.
\end{proof}


\begin{lemma}
\label{lemma-limit-essentially-finite-presentation}
Suppose $R \to S$ is a local homomorphism of local rings.
Assume that $S$ is essentially of finite presentation over $R$.
Then there exists a directed set $(\Lambda, \leq)$, and
a system of local homomorphism $R_\lambda \to S_\lambda$
of local rings such that
\begin{enumerate}
\item The colimit of the system $R_\lambda \to S_\lambda$
is equal to $R \to S$.
\item Each $R_\lambda$ is essentially of finite type
over $\mathbf{Z}$.
\item Each $S_\lambda$ is essentially of finite type
over $R_\lambda$.
\item For each $\lambda \leq \mu$ the map
$S_\lambda \otimes_{R_\lambda} R_\mu \to S_\mu$
presents $S_\mu$ as the localization of
$S_\lambda \otimes_{R_\lambda} R_\mu$
at a prime ideal.
\end{enumerate}
\end{lemma}

\begin{proof}
By assumption we may choose an isomorphism
$\Phi : (R[x_1, \ldots, x_n]/I)_{\mathfrak q} \to S$
where $I \subset R[x_1, \ldots, x_n]$ is a finitely generated ideal,
and $\mathfrak q \subset R[x_1, \ldots, x_n]/I$ is a prime.
(Note that $R \cap \mathfrak q$
is equal to the maximal ideal $\mathfrak m$ of $R$.)
We also choose generators $f_1, \ldots, f_m \in I$ for the ideal $I$.
Write $R$ in any way as a colimit $R = \colim R_\lambda$
over a directed set $(\Lambda, \leq )$, with each $R_\lambda$
local and essentially of finite type over $\mathbf{Z}$.
There exists some $\lambda_0 \in \Lambda$ such that $f_j$ is the image
of some $f_{j, \lambda_0} \in R_{\lambda_0}[x_1, \ldots, x_n]$.
For all $\lambda \geq \lambda_0$ denote
$f_{j, \lambda} \in R_{\lambda}[x_1, \ldots, x_n]$ the image
of $f_{j, \lambda_0}$. Thus we obtain a system of ring maps
$$
R_\lambda[x_1, \ldots, x_n]/(f_{1, \lambda}, \ldots, f_{m, \lambda})
\to
R[x_1, \ldots, x_n]/(f_1, \ldots, f_m) \to S
$$
Set $\mathfrak q_\lambda$ the inverse image of $\mathfrak q$.
Set $S_\lambda = (R_\lambda[x_1, \ldots, x_n]/
(f_{1, \lambda}, \ldots, f_{m, \lambda}))_{\mathfrak q_\lambda}$.
We leave it to the reader to see that this works.
\end{proof}


\begin{lemma}
\label{lemma-limit-module-essentially-finite-presentation}
Suppose $R \to S$ is a local homomorphism of local rings.
Assume that $S$ is essentially of finite presentation over $R$.
Let $M$ be a finitely presented $S$-module.
Then there exists a directed set $(\Lambda, \leq)$, and
a system of local homomorphisms $R_\lambda \to S_\lambda$
of local rings together with $S_\lambda$-modules $M_\lambda$,
such that
\begin{enumerate}
\item The colimit of the system $R_\lambda \to S_\lambda$
is equal to $R \to S$. The colimit of the system $M_\lambda$
is $M$.
\item Each $R_\lambda$ is essentially of finite type
over $\mathbf{Z}$.
\item Each $S_\lambda$ is essentially of finite type
over $R_\lambda$.
\item Each $M_\lambda$ is finite over $S_\lambda$.
\item For each $\lambda \leq \mu$ the map
$S_\lambda \otimes_{R_\lambda} R_\mu \to S_\mu$
presents $S_\mu$ as the localization of
$S_\lambda \otimes_{R_\lambda} R_\mu$
at a prime ideal.
\item For each $\lambda \leq \mu$ the map
$M_\lambda \otimes_{S_\lambda} S_\mu \to M_\mu$
is an isomorphism.
\end{enumerate}
\end{lemma}

\begin{proof}
As in the proof of Lemma \ref{lemma-limit-essentially-finite-presentation}
we may first write $R = \colim R_\lambda$ as a directed colimit
of local $\mathbf{Z}$-algebras which are essentially of finite type.
Next, we may assume that for some $\lambda_1 \in \Lambda$ there
exist $f_{j, \lambda_1} \in R_{\lambda_1}[x_1, \ldots, x_n]$
such that
$$
S =
\colim_{\lambda \geq \lambda_1} S_\lambda, \text{ with }
S_\lambda =
(R_\lambda[x_1, \ldots, x_n]/
(f_{1, \lambda}, \ldots, f_{m, \lambda}))_{\mathfrak q_\lambda}
$$
Choose a presentation
$$
S^{\oplus s} \to S^{\oplus t} \to M \to 0
$$
of $M$ over $S$. Let $A \in \text{Mat}(t \times s, S)$ be
the matrix of the presentation. For some $\lambda_2 \in \Lambda$,
$\lambda_2 \geq \lambda_1$
we can find a matrix $A_{\lambda_2} \in \text{Mat}(t \times s, S_{\lambda_2})$
which maps to $A$. For all $\lambda \geq \lambda_2$ we let
$M_\lambda = \Coker(S_\lambda^{\oplus s} \xrightarrow{A_\lambda}
S_\lambda^{\oplus t})$. We leave it to the reader to see that
this works.
\end{proof}


\begin{lemma}
\label{lemma-colimit-eventually-flat}
Let $R \to S$, $M$, $\Lambda$, $R_\lambda \to S_\lambda$, $M_\lambda$
be as in Lemma \ref{lemma-limit-module-essentially-finite-presentation}.
Assume that $M$ is flat over $R$.
Then for some $\lambda \in \Lambda$ the module
$M_\lambda$ is flat over $R_\lambda$.
\end{lemma}

\begin{proof}
Pick some $\lambda \in \Lambda$ and consider
$$
\text{Tor}_1^{R_\lambda}(M_\lambda, R_\lambda/\mathfrak m_\lambda)
=
\Ker(\mathfrak m_\lambda \otimes_{R_\lambda} M_\lambda
\to M_\lambda).
$$
See Remark \href{https://stacks.math.columbia.edu/tag/00M6}{remark Tor ring mod ideal (Tag 00M6)}. The right hand side
shows that this is a finitely generated $S_\lambda$-module (because
$S_\lambda$ is Noetherian and the modules in question are finite).
Let $\xi_1, \ldots, \xi_n$ be generators.
Because $M$ is flat over $R$ we
have that $0 = \Ker(\mathfrak m_\lambda R \otimes_R M \to M)$.
Since $\otimes$ commutes with colimits we see there exists
a $\lambda' \geq \lambda$ such that each $\xi_i$ maps to
zero in
$\mathfrak m_{\lambda}R_{\lambda'} \otimes_{R_{\lambda'}} M_{\lambda'}$.
Hence we see that
$$
\text{Tor}_1^{R_\lambda}(M_\lambda, R_\lambda/\mathfrak m_\lambda)
\longrightarrow
\text{Tor}_1^{R_{\lambda'}}(M_{\lambda'},
R_{\lambda'}/\mathfrak m_{\lambda}R_{\lambda'})
$$
is zero. Note that
$M_\lambda \otimes_{R_\lambda} R_\lambda/\mathfrak m_\lambda$
is flat over $R_\lambda/\mathfrak m_\lambda$ because this last
ring is a field. Hence we may apply Lemma
\ref{lemma-another-variant-local-criterion-flatness}
to get that $M_{\lambda'}$ is flat over $R_{\lambda'}$.
\end{proof}


\begin{lemma}[Crit\`ere de platitude par fibres; Noetherian case]
\label{lemma-criterion-flatness-fibre-Noetherian}
Let $R$, $S$, $S'$ be Noetherian local rings and let $R \to S \to S'$
be local ring homomorphisms. Let $\mathfrak m \subset R$ be the
maximal ideal. Let $M$ be an $S'$-module. Assume
\begin{enumerate}
\item The module $M$ is finite over $S'$.
\item The module $M$ is not zero.
\item The module $M/\mathfrak m M$
is a flat $S/\mathfrak m S$-module.
\item The module $M$ is a flat $R$-module.
\end{enumerate}
Then $S$ is flat over $R$ and $M$ is a flat $S$-module.
\end{lemma}

\begin{proof}
Set $I = \mathfrak mS \subset S$. Then we see that $M/IM$ is a flat
$S/I$-module because of (3). Since
$\mathfrak m \otimes_R S' \to I \otimes_S S'$ is surjective we see
that also $\mathfrak m \otimes_R M \to I \otimes_S M$ is surjective.
Consider
$$
\mathfrak m \otimes_R M \to I \otimes_S M \to M.
$$
As $M$ is flat over $R$ the composition is injective
and so both arrows are injective.
In particular $\text{Tor}_1^S(S/I, M) = 0$ see
Remark \href{https://stacks.math.columbia.edu/tag/00M6}{remark Tor ring mod ideal (Tag 00M6)}. By
Lemma \ref{lemma-variant-local-criterion-flatness} we conclude
that $M$ is flat over $S$. Note that since $M/\mathfrak m_{S'}M$
is not zero by Nakayama's Lemma \href{https://stacks.math.columbia.edu/tag/00DV}{lemma NAK (Tag 00DV)}
we see that actually $M$ is faithfully flat over $S$ by
Lemma \href{https://stacks.math.columbia.edu/tag/00HP}{lemma ff (Tag 00HP)} (since it forces $M/\mathfrak m_SM \not = 0$).

\medskip\noindent
Consider the exact sequence
$0 \to \mathfrak m \to R \to \kappa \to 0$.
This gives an exact sequence
$0 \to \text{Tor}_1^R(\kappa, S) \to \mathfrak m \otimes_R S \to I \to 0$.
Since $M$ is flat over $S$ this gives an exact sequence
$0 \to \text{Tor}_1^R(\kappa, S)\otimes_S M \to
\mathfrak m \otimes_R M \to I \otimes_S M \to 0$.
By the above this implies that $\text{Tor}_1^R(\kappa, S)\otimes_S M = 0$.
Since $M$ is faithfully flat over $S$ this implies that
$\text{Tor}_1^R(\kappa, S) = 0$ and we conclude that
$S$ is flat over $R$ by Lemma \ref{lemma-local-criterion-flatness}.
\end{proof}


\begin{lemma}
\label{lemma-another-variant-local-criterion-flatness}
Let
$$
\xymatrix{
S \ar[r] & S' \\
R \ar[r] \ar[u] & R' \ar[u]
}
$$
be a commutative diagram of local homomorphisms of local Noetherian rings.
Let $I \subset R$ be a proper ideal.
Let $M$ be a finite $S$-module.
Denote $I' = IR'$ and $M' = M \otimes_S S'$.
Assume that
\begin{enumerate}
\item $S'$ is a localization of the tensor product
$S \otimes_R R'$,
\item $M/IM$ is flat over $R/I$,
\item $\text{Tor}_1^R(M, R/I) \to \text{Tor}_1^{R'}(M', R'/I')$
is zero.
\end{enumerate}
Then $M'$ is flat over $R'$.
\end{lemma}

\begin{proof}
Since $S'$ is a localization of $S \otimes_R R'$ we see that
$M'$ is a localization of $M \otimes_R R'$. Note that
by Lemma \href{https://stacks.math.columbia.edu/tag/00HI}{lemma flat base change (Tag 00HI)} the module $M/IM \otimes_{R/I} R'/I'
= M \otimes_R R' /I'(M \otimes_R R')$ is flat over $R'/I'$. Hence also
$M'/I'M'$ is flat over $R'/I'$ as the localization of a flat module
is flat. By Lemma \ref{lemma-variant-local-criterion-flatness}
it suffices to show that $\text{Tor}_1^{R'}(M', R'/I')$ is zero.
Since $M'$ is a localization of $M \otimes_R R'$, the last assumption
implies that it suffices to show that
$\text{Tor}_1^R(M, R/I) \otimes_R R'
\to
\text{Tor}_1^{R'}(M \otimes_R R', R'/I')$
is surjective.

\medskip\noindent
By Lemma \ref{lemma-surjective-on-tor-one-trivial} we see that
$\text{Tor}_1^R(M, R'/I') \to \text{Tor}_1^{R'}(M \otimes_R R', R'/I')$
is surjective. So now it suffices to show that
$\text{Tor}_1^R(M, R/I) \otimes_R R'
\to
\text{Tor}_1^R(M, R'/I')$
is surjective. This follows from Lemma \ref{lemma-surjective-on-tor-one}
by looking at the ring maps $R \to R/I \to R'/I'$ and the module $M$.
\end{proof}


\begin{lemma}[Variant of the local criterion]
\label{lemma-variant-local-criterion-flatness}
Let $R \to S$ be a local homomorphism of Noetherian
local rings. Let $I \not = R$ be an ideal in $R$.
Let $M$ be a finite $S$-module. If $\text{Tor}_1^R(M, R/I) = 0$
and $M/IM$ is flat over $R/I$, then $M$ is flat over $R$.
\end{lemma}

\begin{proof}
First proof: By
Lemma \href{https://stacks.math.columbia.edu/tag/051C}{lemma what does it mean (Tag 051C)}
we see that $\text{Tor}_1^R(\kappa, M)$ is zero where $\kappa$
is the residue field of $R$. Hence we see that $M$
is flat over $R$ by
Lemma \ref{lemma-local-criterion-flatness}.

\medskip\noindent
Second proof: Let $\mathfrak m$ be the maximal ideal of $R$.
We will show that $\mathfrak m \otimes_R M \to M$ is injective,
and then apply
Lemma \ref{lemma-local-criterion-flatness}.
Suppose that $\sum f_i \otimes x_i \in \mathfrak m \otimes_R M$
and that $\sum f_i x_i = 0$ in $M$. By the equational criterion
for flatness Lemma \href{https://stacks.math.columbia.edu/tag/00HK}{lemma flat eq (Tag 00HK)} applied to $M/IM$
over $R/I$ we see there exist $\overline{a}_{ij} \in R/I$
and $\overline{y}_j \in M/IM$ such that
$x_i \bmod IM = \sum_j \overline{a}_{ij} \overline{y}_j $
and $0 = \sum_i (f_i \bmod I) \overline{a}_{ij}$.
Let $a_{ij} \in R$ be a lift of $\overline{a}_{ij}$ and
similarly let $y_j \in M$ be a lift of $\overline{y}_j$.
Then we see that
\begin{eqnarray*}
\sum f_i \otimes x_i
& = &
\sum f_i \otimes x_i +
\sum f_ia_{ij} \otimes y_j -
\sum f_i \otimes a_{ij} y_j
\\
& = &
\sum f_i \otimes (x_i - \sum a_{ij} y_j) +
\sum (\sum f_i a_{ij}) \otimes y_j
\end{eqnarray*}
Since $x_i - \sum a_{ij} y_j \in IM$ and
$\sum f_i a_{ij} \in I$ we see that there exists
an element in $I \otimes_R M$ which maps to our given
element $\sum f_i \otimes x_i$ in $\mathfrak m \otimes_R M$.
But $I \otimes_R M \to M$ is injective by assumption (see
Remark \href{https://stacks.math.columbia.edu/tag/00M6}{remark Tor ring mod ideal (Tag 00M6)}) and we win.
\end{proof}


\begin{lemma}[Local criterion for flatness]
\label{lemma-local-criterion-flatness}
Let $R \to S$ be a local homomorphism of local Noetherian
rings. Let $\mathfrak m$ be the maximal ideal of $R$,
and let $\kappa = R/\mathfrak m$.
Let $M$ be a finite $S$-module. If $\text{Tor}_1^R(\kappa, M) = 0$,
then $M$ is flat over $R$.
\end{lemma}

\begin{proof}
Let $I \subset R$ be an ideal. By Lemma \href{https://stacks.math.columbia.edu/tag/00HD}{lemma flat (Tag 00HD)} it suffices
to show that $I \otimes_R M \to M$ is injective. By Remark
\href{https://stacks.math.columbia.edu/tag/00M6}{remark Tor ring mod ideal (Tag 00M6)} we see that this kernel is
equal to $\text{Tor}_1^R(M, R/I)$. By
Lemma \ref{lemma-prepare-local-criterion-flatness}
we see that $J \otimes_R M \to M$ is injective for all ideals
of finite colength.

\medskip\noindent
Choose $n >> 0$ and consider the following short exact
sequence
$$
0
\to I \cap \mathfrak m^n
\to I \oplus \mathfrak m^n
\to I + \mathfrak m^n
\to 0
$$
This is a sub sequence of the short exact sequence
$0 \to R \to R^{\oplus 2} \to R \to 0$. Thus we get the diagram
$$
\xymatrix{
(I\cap \mathfrak m^n) \otimes_R M \ar[r] \ar[d] &
I \otimes_R M \oplus \mathfrak m^n \otimes_R M \ar[r] \ar[d] &
(I + \mathfrak m^n) \otimes_R M \ar[d] \\
M \ar[r] &
M \oplus M \ar[r] &
M
}
$$
Note that $I + \mathfrak m^n$ and $\mathfrak m^n$
are ideals of finite colength.
Thus a diagram chase shows that
$\Ker((I \cap \mathfrak m^n)\otimes_R M \to M)
\to \Ker(I \otimes_R M \to M)$
is surjective. We conclude in particular that
$K = \Ker(I \otimes_R M \to M)$ is contained
in the image of $(I \cap \mathfrak m^n) \otimes_R M$
in $I \otimes_R M$. By Artin-Rees, Lemma \href{https://stacks.math.columbia.edu/tag/00IN}{lemma Artin Rees (Tag 00IN)}
we see that $K$ is contained
in $\mathfrak m^{n-c}(I \otimes_R M)$ for some $c > 0$
and all $n >> 0$. Since $I \otimes_R M$ is a finite
$S$-module (!) and since $S$ is Noetherian, we see
that this implies $K = 0$. Namely, the above implies
$K$ maps to zero in the $\mathfrak mS$-adic completion
of $I \otimes_R M$. But the map from $S$
to its $\mathfrak mS$-adic completion is faithfully flat
by Lemma \href{https://stacks.math.columbia.edu/tag/00MC}{lemma completion faithfully flat (Tag 00MC)}.
Hence $K = 0$, as desired.
\end{proof}


\begin{lemma}
\label{lemma-surjective-on-tor-one}
Let $R \to R' \to R''$ be ring maps.
Let $M$ be an $R$-module. Suppose that $M \otimes_R R'$
is flat over $R'$. Then the natural map
$\text{Tor}_1^R(M, R') \otimes_{R'} R'' \to
\text{Tor}_1^R(M, R'')$ is onto.
\end{lemma}

\begin{proof}
Let $F_\bullet$ be a free resolution of $M$ over $R$.
The complex $F_2 \otimes_R R' \to F_1\otimes_R R' \to F_0 \otimes_R R'$
computes $\text{Tor}_1^R(M, R')$.
The complex $F_2 \otimes_R R'' \to F_1\otimes_R R'' \to F_0 \otimes_R R''$
computes $\text{Tor}_1^R(M, R'')$. Note that
$F_i \otimes_R R' \otimes_{R'} R'' = F_i \otimes_R R''$. Let
$K' = \Ker(F_1\otimes_R R' \to F_0 \otimes_R R')$ and
similarly $K'' = \Ker(F_1\otimes_R R'' \to F_0 \otimes_R R'')$.
Thus we have an exact sequence
$$
0 \to K' \to F_1\otimes_R R' \to F_0 \otimes_R R' \to M \otimes_R R' \to 0.
$$
By the assumption that $M \otimes_R R'$ is flat over $R'$,
the sequence
$$
K' \otimes_{R'} R'' \to
F_1 \otimes_R R'' \to
F_0 \otimes_R R'' \to
M \otimes_R R'' \to 0
$$
is still exact. This means that $K' \otimes_{R'} R'' \to K''$
is surjective. Since $\text{Tor}_1^R(M, R')$ is a quotient of $K'$ and
$\text{Tor}_1^R(M, R'')$ is a quotient of $K''$ we win.
\end{proof}


\begin{lemma}
\label{lemma-surjective-on-tor-one-trivial}
Let $R \to R'$ be a ring map. Let $I \subset R$ be
an ideal and $I' = IR'$. Let $M$ be an $R$-module
and set $M' = M \otimes_R R'$. The natural map
$\text{Tor}_1^R(R'/I', M) \to \text{Tor}_1^{R'}(R'/I', M')$
is surjective.
\end{lemma}

\begin{proof}
Let $F_2 \to F_1 \to F_0 \to M \to 0$ be a free resolution of
$M$ over $R$. Set $F_i' = F_i \otimes_R R'$. The sequence
$F_2' \to F_1' \to F_0' \to M' \to 0$ may no longer be exact
at $F_1'$. A free resolution of $M'$ over $R'$ therefore looks
like
$$
F_2' \oplus F_2'' \to F_1' \to F_0' \to M' \to 0
$$
for a suitable free module $F_2''$ over $R'$. Next, note that
$F_i \otimes_R R'/I' = F_i' / IF_i' = F_i'/I'F_i'$.
So the complex $F_2'/I'F_2' \to F_1'/I'F_1' \to F_0'/I'F_0'$
computes $\text{Tor}_1^R(M, R'/I')$. On the other hand
$F_i' \otimes_{R'} R'/I' = F_i'/I'F_i'$ and similarly
for $F_2''$. Thus the complex
$F_2'/I'F_2' \oplus F_2''/I'F_2'' \to F_1'/I'F_1' \to F_0'/I'F_0'$
computes $\text{Tor}_1^{R'}(M', R'/I')$. Since the vertical
map on complexes
$$
\xymatrix{
F_2'/I'F_2' \ar[r] \ar[d] &
F_1'/I'F_1' \ar[r] \ar[d] &
F_0'/I'F_0' \ar[d] \\
F_2'/I'F_2' \oplus F_2''/I'F_2'' \ar[r] &
F_1'/I'F_1' \ar[r] &
F_0'/I'F_0'
}
$$
clearly induces a surjection on cohomology we win.
\end{proof}


\begin{lemma}
\label{lemma-prepare-local-criterion-flatness}
Let $R$ be a local ring with maximal ideal $\mathfrak m$
and residue field $\kappa = R/\mathfrak m$.
Let $M$ be an $R$-module. If $\text{Tor}_1^R(\kappa, M) = 0$,
then for every finite length $R$-module $N$ we have
$\text{Tor}_1^R(N, M) = 0$.
\end{lemma}

\begin{proof}
By descending induction on the length of $N$.
If the length of $N$ is $1$, then $N \cong \kappa$
and we are done. If the length of $N$ is more than
$1$, then we can fit $N$ into a short exact sequence
$0 \to N' \to N \to N'' \to 0$ where $N'$, $N''$ are
finite length $R$-modules of smaller length.
The vanishing of $\text{Tor}_1^R(N, M)$ follows
from the vanishing of $\text{Tor}_1^R(N', M)$
and $\text{Tor}_1^R(N'', M)$ (induction hypothesis)
and the long exact sequence of Tor groups, see Lemma
\href{https://stacks.math.columbia.edu/tag/00M0}{lemma long exact sequence tor (Tag 00M0)}.
\end{proof}


\begin{lemma}[Crit\`ere de platitude par fibres]
\label{lemma-criterion-flatness-fibre}
Let $R$, $S$, $S'$ be local rings and let $R \to S \to S'$ be local ring
homomorphisms. Let $M$ be an $S'$-module. Let $\mathfrak m \subset R$
be the maximal ideal. Assume
\begin{enumerate}
\item The ring maps $R \to S$ and $R \to S'$ are essentially
of finite presentation.
\item The module $M$ is of finite presentation over $S'$.
\item The module $M$ is not zero.
\item The module $M/\mathfrak mM$ is a flat $S/\mathfrak mS$-module.
\item The module $M$ is a flat $R$-module.
\end{enumerate}
Then $S$ is flat over $R$ and $M$ is a flat $S$-module.
\end{lemma}

\begin{proof}
As in the proof of Lemma \ref{lemma-limit-essentially-finite-presentation}
we may first write $R = \colim R_\lambda$ as a directed colimit
of local $\mathbf{Z}$-algebras which are essentially of finite type.
Denote $\mathfrak p_\lambda$ the maximal ideal of $R_\lambda$.
Next, we may assume that for some $\lambda_1 \in \Lambda$ there
exist $f_{j, \lambda_1} \in R_{\lambda_1}[x_1, \ldots, x_n]$
such that
$$
S =
\colim_{\lambda \geq \lambda_1} S_\lambda, \text{ with }
S_\lambda =
(R_\lambda[x_1, \ldots, x_n]/
(f_{1, \lambda}, \ldots, f_{u, \lambda}))_{\mathfrak q_\lambda}
$$
For some $\lambda_2 \in \Lambda$,
$\lambda_2 \geq \lambda_1$ there exist
$g_{j, \lambda_2} \in R_{\lambda_2}[x_1, \ldots, x_n, y_1, \ldots, y_m]$
with images
$\overline{g}_{j, \lambda_2} \in S_{\lambda_2}[y_1, \ldots, y_m]$
such that
$$
S' =
\colim_{\lambda \geq \lambda_2} S'_\lambda, \text{ with }
S'_\lambda =
(S_\lambda[y_1, \ldots, y_m]/
(\overline{g}_{1, \lambda}, \ldots,
\overline{g}_{v, \lambda}))_{\overline{\mathfrak q}'_\lambda}
$$
Note that this also implies that
$$
S'_\lambda =
(R_\lambda[x_1, \ldots, x_n, y_1, \ldots, y_m]/
(g_{1, \lambda}, \ldots, g_{v, \lambda}))_{\mathfrak q'_\lambda}
$$
Choose a presentation
$$
(S')^{\oplus s} \to (S')^{\oplus t} \to M \to 0
$$
of $M$ over $S'$. Let $A \in \text{Mat}(t \times s, S')$ be
the matrix of the presentation. For some $\lambda_3 \in \Lambda$,
$\lambda_3 \geq \lambda_2$
we can find a matrix $A_{\lambda_3} \in \text{Mat}(t \times s, S_{\lambda_3})$
which maps to $A$. For all $\lambda \geq \lambda_3$ we let
$M_\lambda = \Coker((S'_\lambda)^{\oplus s} \xrightarrow{A_\lambda}
(S'_\lambda)^{\oplus t})$.

\medskip\noindent
With these choices, we have for each $\lambda_3 \leq \lambda \leq \mu$
that $S_\lambda \otimes_{R_{\lambda}} R_\mu \to S_\mu$ is a localization,
$S'_\lambda \otimes_{S_{\lambda}} S_\mu \to S'_\mu$ is a localization, and
the map $M_\lambda \otimes_{S'_\lambda} S'_\mu \to M_\mu$ is an
isomorphism. This also implies that
$S'_\lambda \otimes_{R_{\lambda}} R_\mu \to S'_\mu$ is a localization.
Thus, since $M$ is flat over $R$ we see by
Lemma \ref{lemma-colimit-eventually-flat} that
for all $\lambda$ big enough the module $M_\lambda$ is
flat over $R_\lambda$.
Moreover, note that
$
\mathfrak m = \colim \mathfrak p_\lambda
$,
$
S/\mathfrak mS = \colim S_\lambda/\mathfrak p_\lambda S_\lambda
$,
$
S'/\mathfrak mS' = \colim S'_\lambda/\mathfrak p_\lambda S'_\lambda
$,
and
$
M/\mathfrak mM = \colim M_\lambda/\mathfrak p_\lambda M_\lambda
$. Also, for each $\lambda_3 \leq \lambda \leq \mu$ we see (from the
properties listed above) that
$$
S'_\lambda/\mathfrak p_\lambda S'_\lambda
\otimes_{S_{\lambda}/\mathfrak p_\lambda S_\lambda}
S_\mu/\mathfrak p_\mu S_\mu
\longrightarrow
S'_\mu/\mathfrak p_\mu S'_\mu
$$
is a localization, and the map
$$
M_\lambda / \mathfrak p_\lambda M_\lambda
\otimes_{S'_\lambda/\mathfrak p_\lambda S'_\lambda}
S'_\mu /\mathfrak p_\mu S'_\mu
\longrightarrow
M_\mu/\mathfrak p_\mu M_\mu
$$
is an isomorphism. Hence the system
$(S_\lambda/\mathfrak p_\lambda S_\lambda \to
S'_\lambda/\mathfrak p_\lambda S'_\lambda,
M_\lambda/\mathfrak p_\lambda M_\lambda)$
is a system as in
Lemma \ref{lemma-limit-module-essentially-finite-presentation} as well.
We may apply Lemma \ref{lemma-colimit-eventually-flat} again because
$M/\mathfrak m M$ is assumed flat over $S/\mathfrak mS$ and we see that
$M_\lambda/\mathfrak p_\lambda M_\lambda$ is flat over
$S_\lambda/\mathfrak p_\lambda S_\lambda$ for all $\lambda$ big enough.
Thus for $\lambda$ big enough the data
$R_\lambda \to S_\lambda \to S'_\lambda, M_\lambda$ satisfies
the hypotheses of Lemma \ref{lemma-criterion-flatness-fibre-Noetherian}.
Pick such a $\lambda$. Then $S = S_\lambda \otimes_{R_\lambda} R$
is flat over $R$, and $M = M_\lambda \otimes_{S_\lambda} S$
is flat over $S$ (since the base change of a flat module is flat).
\end{proof}
\end{document}
